
\chapter{Diagrama de Robustez}\label{ch:diagramarobustez}

El diagrama de robustez es un artefacto del método ICONIX que sirve de puente entre el análisis (casos de uso, modelo conceptual) y el diseño detallado. A diferencia del diagrama de clases conceptuales del Modelo de Análisis (que solo muestra conceptos del problema), aquí ya se modelan \textbf{clases de software}, clasificadas en tres categorías:

\begin{itemize}
  \item \textbf{Boundary (frontera):} interfaces de usuario y puntos de integración. En FisioManager: pantallas (\codeinline{UIAgendaPaciente}, \codeinline{UIRegistroFeedback}, \codeinline{UIMensajeria}, \ldots), el servicio de notificaciones (correo), el servicio de tiempo real (\codeinline{ServicioTiempoReal}, Socket.IO) y el servicio de IA (\codeinline{ServicioIA}: Whisper para transcribir y Ollama para resumir).
  \item \textbf{Control:} lógica de aplicación que orquesta el flujo entre fronteras y entidades. En FisioManager: \codeinline{GestorCitas}, \codeinline{GestorFeedback}, \codeinline{AnalizadorIA}, \codeinline{GestorMensajeria}, \codeinline{GestorNotificaciones}, \codeinline{GestorPlanTratamiento}, \codeinline{GestorNotasClinicas}, etc.
  \item \textbf{Entity (entidad):} entidades persistentes del dominio, generalmente derivadas del diagrama de clases conceptuales.
\end{itemize}

\section{Flujos consolidados}

Para mantener una visión integrada del sistema se presenta un único diagrama de robustez consolidado que muestra seis flujos representativos:

\begin{enumerate}
  \item \textbf{Agendar cita (Paciente):} el paciente consulta franjas libres y solicita una cita; el \codeinline{GestorCitas} delega en \codeinline{ValidadorDisponibilidad} para verificar horario y colisiones, crea la \codeinline{Cita}, avisa a la \codeinline{ListaEspera} si procede y, vía \codeinline{GestorNotificaciones}, genera la \codeinline{Notificacion} (correo y aviso en vivo).
  \item \textbf{Registrar feedback de sesión (Paciente):} el paciente envía audio, dolor, emoción y mapa de dolor; el \codeinline{GestorFeedback} crea el \codeinline{Feedback} y los \codeinline{PuntoDolor}, y delega en \codeinline{AnalizadorIA}, que usa \codeinline{ServicioIA} (Whisper + Ollama) y escribe la transcripción y el resumen \emph{en el propio} \codeinline{Feedback}.
  \item \textbf{Mensajería en tiempo real:} el paciente escribe un mensaje; el \codeinline{GestorMensajeria} ubica/crea la \codeinline{Conversacion}, persiste el \codeinline{Mensaje} y lo entrega al terapeuta mediante \codeinline{ServicioTiempoReal} (Socket.IO), notificando al receptor.
  \item \textbf{Elaborar plan de tratamiento (Terapeuta):} el terapeuta crea o edita un \codeinline{PlanTratamiento} con sus \codeinline{FaseTratamiento} y prescribe \codeinline{AsignacionRutina}.
  \item \textbf{Registrar nota clínica (Terapeuta):} desde el expediente, el terapeuta agrega una \codeinline{NotaClinica} que el \codeinline{GestorNotasClinicas} adjunta al \codeinline{HistorialClinico}.
  \item \textbf{Consulta de agenda (Terapeuta):} el terapeuta consulta su agenda del día; el \codeinline{GestorCitas} devuelve las citas programadas.
\end{enumerate}

\begin{figure}[H]
  \centering
  \includegraphics[width=\textwidth,keepaspectratio]{FisioManager_DiagramaRobustez.png}
  \caption{Diagrama de Robustez consolidado: clases de software (Boundary, Control, Entity) que participan en los flujos principales de FisioManager.}
  \label{fig:diagrama-robustez}
\end{figure}

\section{Reglas de robustez aplicadas}

Se respetan las reglas clásicas de ICONIX:

\begin{itemize}
  \item Los actores solo interactúan con clases \textbf{boundary}.
  \item Las clases \textbf{boundary} no se comunican entre sí directamente: siempre pasan por una clase \textbf{control}.
  \item Las clases \textbf{entity} no se comunican entre sí ni con \textbf{boundary} directamente: la coordinación entre entidades es responsabilidad de las clases \textbf{control}.
  \item Las clases \textbf{control} pueden hablar con cualquier otro tipo de clase.
\end{itemize}

Este diagrama es la base para identificar los \textit{servicios de aplicación} y \textit{repositorios} que aparecerán en el diseño detallado, y para construir más adelante los diagramas de secuencia por caso de uso.

\section{Código PlantUML del diagrama}

El código fuente se entrega como \codeinline{08\_diagrama\_robustez.puml} en el directorio \codeinline{puml/}. Se regenera con:

\begin{lstlisting}[language=bash]
plantuml -tpng 08_diagrama_robustez.puml
\end{lstlisting}

El código fuente completo se incluye en el anexo de código fuente PlantUML del Modelo de Análisis.
